% ============ 6 平展拓扑的层 (原书 pp.42-49) ============
\chapter{平展拓扑的层}\label{sec:6}
\index{Galois 覆盖}
\index{一般 Galois}
\index{层}
\index{常值层}
\index{摩天楼层}
\index{局部常值层}
\index{茎}
\index{忠实平坦}
\index{结构层}

设 $X$ 是簇(或概形).按一般定义，$X_{\mathrm{et}}$ 上的层 $\mathcal{F}$ 是反变函子 $\mathbf{Et}/X \to \mathbf{Sets}$(或 $\mathbf{Ab}$，或 $\ldots$)，使得对每个平展映射 $U \to X$ 与每个平展覆盖 $(U_i \to U)$，
\[
(\mathrm{S}) \colon \quad \mathcal{F}(U) \to \prod_{i \in I} \mathcal{F}(U_i) \rightrightarrows \prod_{(i, j) \in I \times I} \mathcal{F}(U_i \times_{U} U_j)
\]
正合.

注意 $X_{\mathrm{et}}$ 上的层 $\mathcal{F}$ 通过限制对每个平展映射 $U \to X$ 定义了 $U_{\mathrm{zar}}$ 上的层.特别地，若 $U = \coprod U_i$，则 $\mathcal{F}(U) \xrightarrow{\sim} \prod \mathcal{F}(U_i)$.把层条件应用于 $I$ 为空集的情形，发现 $\mathcal{F}(\varnothing)$ 是单元素集合(若 $\mathcal{F}$ 是群层则为零群).这里 $\varnothing$ 是空簇(或空概形).

对 Zariski 覆盖，层条件 (S) 就是熟悉的那个.下面我们考察它对 Galois 覆盖意味着什么.

\subsection*{Galois 覆盖}

设 $\varphi \colon Y \to X$ 是态射，$G$ 是有限群.$G$ 在 $Y$ 上(在 $X$ 上)的右作用是映射 $\alpha \colon G \to \operatorname{Aut}_X(Y)$，使得 $\alpha(gh) = \alpha(h) \circ \alpha(g)$.若 $X$ 与 $Y$ 仿射，使 $\varphi$ 对应于环的映射 $A \to B$，则给出 $G$ 在 $Y$ 上(在 $X$ 上)的右作用与给出 $G$ 在 $A$-代数 $B$ 上的左作用是一回事.

\begin{definition}\label{def:6.1}
设 $Y \to X$ 是忠实平坦映射，$G$ 是作用在 $Y$ 上(在 $X$ 上)的有限群.若态射
\[
Y \times G \to Y \times_{X} Y, \qquad (y, g) \mapsto (y, yg)
\]
是同构，则 $Y \to X$ 称为带群 $G$ 的 $X$ 的 Galois 覆盖.
\end{definition}

这里 $Y \times G$ 表示由 $G$ 的元素指标化的 $Y$ 的副本的不交并：
\[
Y \times G = \coprod_{g} Y_g, \qquad Y_g = Y.
\]
映射 $\mathrm{id} \colon Y_g \to Y$ 与 $g \colon Y_g \to Y$ 定义映射 $Y_g \to Y \times_{X} Y$，而定义中的条件要求这些映射诱导同构 $\coprod_{g} Y_g \to Y \times_{X} Y$.

若 $\varphi \colon Y \to X$ 是 Galois 覆盖，则 $\varphi$ 满，有限且平展，次数等于 $G$ 的阶(因为作平坦基变换后变成这样).反过来，可以证明：若 $Y \to X$ 满，有限且平展，次数等于 $\operatorname{Aut}_X(Y)$ 的阶，则它是带群 $\operatorname{Aut}_X(Y)$ 的 Galois 覆盖.

\begin{remark}\label{rem:6.2}
不可约簇的态射 $\varphi \colon Y \to X$ 称为一般 Galois 的，如果 $k(Y)$ 在 $k(X)$ 上 Galois.有些坏蛋把``一般''去掉，把这样的态射称为 Galois 的.

$A$-代数 $B$ 称为 Galois 的，如果存在作用在 $B$ 上(通过 $A$-代数自同构)的群 $G$，使得 $\operatorname{Spec} B \to \operatorname{Spec} A$ 是带群 $G$ 的 Galois 覆盖.显式地，这意味着 $A \to B$ 忠实平坦，且同态
\[
B \otimes_{A} B \to \prod_{g \in G} B, \qquad b \otimes b' \mapsto (\ldots, b \cdot g(b'), \ldots),
\]
是同构.这里 $\prod_{g \in G} B$ 表示由 $G$ 的元素指标化的 $B$ 的副本的乘积.
\end{remark}

\begin{example}\label{ex:6.3}
设 $K = k[T]/(f(T))$，其中 $k$ 是域，$f(T)$ 是 $k[T]$ 中的首一不可约多项式.若
\[
f(T) = f_1(T)^{e_1} \cdots f_r(T)^{e_r}
\]
是 $f(T)$ 在 $K[T]$ 中分解为互异不可约多项式的幂，则由中国剩余定理，
\[
K \otimes_k K \cong K[T]/(f(T)) \cong \prod K[T]/(f_i(T)^{e_i}).
\]
由此推出，$K$ 是 Galois $k$-代数，当且仅当 $f(T)$ 在 $K[T]$ 中分裂为互异的线性因子，即 $f(T)$ 可分且 $K$ 是它的分裂域.

例如，设 $f(T)$ 是 $n$ 次首一不可约可分多项式.若 $K = k[T]/(f)$ 是 $f$ 的分裂域，则 $K/k$ 是 Galois 的(由标准 Galois 理论)，其 Galois 群 $G$ 是 $f$ 的 $n$ 个根的置换群的 $n$ 阶传递子群.
\end{example}

设 $Y \to X$ 是带群 $G$ 的 Galois 覆盖.则 $G$ 右作用于 $Y$，因此对任意预层 $P$，由于 $P$ 是反变函子，$G$ 左作用于 $P(Y)$.

\begin{proposition}\label{prop:6.4}
设 $Y \to X$ 是带群 $G$ 的 Galois 覆盖，$F$ 是 $X_{\mathrm{et}}$ 上把不交并送到乘积的预层.则 $F$ 对覆盖 $Y \to X$ 满足层条件 (S)，当且仅当限制映射 $F(X) \to F(Y)$ 把 $F(X)$ 等同于 $F(Y)$ 中被 $G$ 固定的元素子集 $F(Y)^{G}$.
\end{proposition}

\begin{proof}
有交换图
\[
\begin{tikzcd}
X \arrow[r, shift left] \arrow[r, shift right] & Y \times_{X} Y \\
X \arrow[u, equal] \arrow[r, shift left] \arrow[r, shift right] & Y \times G \arrow[u, "\sim"']
\end{tikzcd}
\]
其中第一行中的投影映射 $(y, y') \mapsto y$ 与 $(y, y') \mapsto y'$ 分别对应于第二行中的映射 $(y, g) \mapsto y$ 与 $(y, g) \mapsto yg$.把 $F$ 应用于该图，得到交换图
\[
\begin{tikzcd}
F(X) \arrow[r] & F(Y) \arrow[r, shift left] \arrow[r, shift right] & F(Y \times_{X} Y) \\
F(X) \arrow[u, equal] \arrow[r] & F(Y) \arrow[u, "\sim"] \arrow[r, shift left] \arrow[r, shift right] & \prod_{g \in G} F(Y) \arrow[u, "\sim"']
\end{tikzcd}
\]
这里 $\prod_{g \in G} F(Y)$ 是由 $G$ 的元素指标化的 $F(Y)$ 的副本的乘积，而映射 $F(Y) \to \prod_{g \in G} F(Y)$ 分别是
\[
s \mapsto (\ldots, s, \ldots, s, \ldots), \qquad s \mapsto (\ldots, gs, \ldots)
\]
这两个映射在 $s \in F(Y)$ 上一致，当且仅当对所有 $g \in G$，$gs = s$.
\end{proof}

\begin{exercise}\label{exc:6.5}
(a) 设 $F$ 是 $X_{\mathrm{et}}$ 上的 Abel 群层，$Y \to X$ 是带群 $G$ 的 Galois 覆盖.证明复形
\[
F(X) \to F(Y) \to F(Y \times_{X} Y) \to F(Y \times_{X} Y \times_{X} Y) \to \cdots
\]
同构于 $G$ 作用在 $F(Y)$ 上的非齐次上链复形(见 CFT p.62).复形中的每个映射是各投影映射给出的映射的交替和.

(b) 证明 6.3 在把 $k$ 换成局部 Hensel 环时仍然成立.
\end{exercise}

\subsection*{成为层的判别准则}

下一个结果使检验预层是否为层更容易.

\begin{proposition}\label{prop:6.6}
为了验证 $X_{\mathrm{et}}$ 上的预层 $F$ 是层，只需检验 $F$ 对 Zariski 开覆盖和由单个映射组成的平展覆盖 $V \to U$($V$ 与 $U$ 都仿射)满足层条件 (S).
\end{proposition}

\begin{proof}
若 $F$ 对 Zariski 开覆盖满足层条件，则 $F(\coprod U_i) = \prod F(U_i)$.由此推出，对指标集 $I$ 有限的覆盖 $(U_i \to U)_{i \in I}$，序列 (S) 同构于由单个态射 $(\coprod U_i \to U)$ 产生的序列 (S)，因为
\[
\left(\coprod U_i\right) \times_{U} \left(\coprod U_i\right) = \coprod_{(i,j) \in I \times I} \left( U_i \times_{U} U_j \right).
\]
由于仿射簇(或概形)的有限不交并仍是仿射的，命题叙述中的第二个条件蕴含 (S) 对指标集有限且 $U$ 与诸 $U_i$ 仿射的覆盖 $(U_i \to U)_{i \in I}$ 正合.

证明的其余部分只涉及简单的图追图，见 EC II 1.5.
\end{proof}

\begin{remark}\label{rem:6.7}
可以证明，每个有限平展映射 $V \to U$ 都被 Galois 映射 $V' \to U$ 支配，即存在穿过 $V \to U$ 的 Galois 覆盖 $V' \to U$，使得 $V' \to V$ 满射.由此以及命题推出，要验证 $X_{\mathrm{fet}}$(有限平展拓扑---见 \S 5)上的预层 $F$ 是层，只需检验 $F$ 对 Zariski 开覆盖与 Galois 覆盖 $V \to U$ 满足层条件 (S).这些论断对平展拓扑不成立.
\end{remark}

\subsection*{\texorpdfstring{$X_{\mathrm{et}}$ 上层的一些例子}{X\_et 上层的一些例子}}

设 $A \to B$ 是对应于仿射簇(或概形)的满平展态射 $V \to U$ 的环同态.在检验 (\ref{prop:6.6}) 中的第二个条件时，我们通常只用 $A \to B$ 忠实平坦\footnote{回忆平坦同态 $A \to B$ 是忠实平坦的，如果它满足下列等价条件之一：
\begin{enumerate}
  \item 若 $A$-模 $M$ 非零，则 $B \otimes_{A} M$ 非零；
  \item 若 $A$-模的序列 $M' \to M \to M''$ 不正合，则 $B \otimes_{A} M' \to B \otimes_{A} M \to B \otimes_{A} M''$ 也不正合；
  \item 映射 $\operatorname{Spec} B \to \operatorname{Spec} A$ 满射(对仿射 $k$-代数，这等价于 $\operatorname{Specm} B \to \operatorname{Specm} A$ 满射).
\end{enumerate}}这一事实(即不需要用它非分歧).

\textbf{$X_{\mathrm{et}}$ 上的结构层} 对任何平展映射 $U \to X$，定义 $\mathcal{O}_{X_{\mathrm{et}}}(U) = \Gamma(U, \mathcal{O}_U)$.当然，对任何在 $X$ 上平展的 $U$，它在 $U_{\mathrm{zar}}$ 上的限制是层.它是 $X_{\mathrm{et}}$ 上的层这一点由命题~\ref{prop:6.6} 与下一个命题推出.

\begin{proposition}\label{prop:6.8}
对每个忠实平坦同态 $A \to B$，序列
\[
0 \to A \to B \xrightarrow{b \mapsto 1 \otimes b - b \otimes 1} B \otimes_{A} B
\]
正合.
\end{proposition}

\begin{proof}
第 1 步：若 $f \colon A \to B$ 有截面，即满足 $s \circ f = \mathrm{id}$ 的同态 $s \colon B \to A$，则论断成立.

为证明这一点，设 $k \colon B \otimes_{A} B \to B$ 把 $b \otimes b'$ 送到 $b \cdot f(s(b'))$.则
\[
k(1 \otimes b - b \otimes 1) = f(s(b)) - b.
\]
因此，若 $1 \otimes b - b \otimes 1 = 0$，则 $b = f(s(b)) \in f(A)$.

第 2 步：若论断对 $a' \mapsto a' \otimes 1 \colon A' \to A' \otimes_{A} B$ 成立，其中 $A \to A'$ 是某个忠实平坦同态，则它对 $A \to B$ 成立.

$A' \to A' \otimes B$ 的序列由 $A \to B$ 的序列张量 $A'$ 得到.

第 3 步：同态 $b \mapsto b \otimes 1 \colon B \to B \otimes_{A} B$ 有截面，即映射 $b \otimes b' \mapsto bb'$.

由于按假设 $A \to B$ 忠实平坦，这就完成了证明.
\end{proof}

\begin{remark}\label{rem:6.9}
设 $K \to L$ 是按上述意义带群 $G$ 的 Galois 域扩张.则 $L^{G} = K$(由命题~\ref{prop:6.4} 与~\ref{prop:6.8})，从而 $L$ 按通常意义在 $K$ 上 Galois(FT 3.9).
\end{remark}

\textbf{由概形 $Z$ 定义的层} $X$-概形 $Z$ 定义反变函子：
\[
F_Z \colon \mathbf{Et}/X \to \mathbf{Sets}, \qquad F_Z(U) = \operatorname{Hom}_X(U, Z).
\]
我断言这是集合层.容易看出 $F_Z$ 对 Zariski 开覆盖满足层判别准则.因此只需证明，对任何忠实平坦映射 $A \to B$，
\[
Z(A) \to Z(B) \rightrightarrows Z(B \otimes_{A} B)
\]
正合.若 $Z$ 仿射，设由环 $C$ 定义，则序列变为
\[
\operatorname{Hom}_{A\text{-alg}}(C, A) \to \operatorname{Hom}_{A\text{-alg}}(C, B) \rightrightarrows \operatorname{Hom}_{A\text{-alg}}(C, B \otimes_{A} B).
\]
它的正合性立即由命题~\ref{prop:6.8} 推出.非仿射 $Z$ 的情形留给读者.

若 $Z$ 有群结构，则 $F_Z$ 是群层.

\begin{example}\label{ex:6.10}
(a) 设 $\mu_n$ 是由单个方程
\[
T^{n} - 1 = 0
\]
定义的簇(或概形).则 $\mu_n(U)$ 是 $\Gamma(U, \mathcal{O}_U)$ 中 $1$ 的 $n$ 次单位根群.

(b) 设 $\mathbb{G}_{a}$ 是把仿射线看作加法群.则 $\mathbb{G}_{a}(U) = \Gamma(U, \mathcal{O}_U)$(看作 Abel 群).

(c) 设 $\mathbb{G}_{m}$ 是去掉原点的仿射线，看作乘法群.则 $\mathbb{G}_{m}(U) = \Gamma(U, \mathcal{O}_U)^{\times}$.

(d) 设 $\mathrm{GL}_n$ 是在 $n^{2} + 1$ 个变量 $T, T_{11}, \ldots, T_{nn}$ 中由单个方程
\[
T \cdot \det(T_{ij}) = 1
\]
定义的簇(或概形).则 $\mathrm{GL}_n(U) = \mathrm{GL}_n(\Gamma(U, \mathcal{O}_U))$，即元素来自环 $\Gamma(U, \mathcal{O}_U)$ 的可逆 $n \times n$ 矩阵群.例如，$\mathrm{GL}_1 = \mathbb{G}_{m}$.
\end{example}

\textbf{常值层} 设 $X$ 是簇或拟紧概形.对每个集合 $\Lambda$，定义
\[
F_{\Lambda}(U) = \Lambda^{\pi_0(U)}
\]
---由 $U$ 的连通分支集 $\pi_0(U)$ 指标化的 $\Lambda$ 的副本的乘积.配上显然的限制映射，这是层，称为由 $\Lambda$ 定义的 $X_{\mathrm{et}}$ 上的常值层.若 $\Lambda$ 有限，则它是由概形 $X \times \Lambda$($X$ 的由 $\Lambda$ 指标化的副本的不交并)定义的层.当 $\Lambda$ 是群时，$F_{\Lambda}$ 是群层.

\textbf{由凝聚 $\mathcal{O}_X$-模定义的层} 设 $\mathcal{M}$ 是通常代数几何意义下 $X_{\mathrm{zar}}$ 上的凝聚 $\mathcal{O}_X$-模层.对每个平展映射 $\varphi \colon U \to X$，得到 $U_{\mathrm{zar}}$ 上的凝聚 $\mathcal{O}_U$-模层 $\varphi^{*}\mathcal{M}$.例如，若 $U$ 与 $X$ 仿射，分别对应环 $B$ 与 $A$，则 $\mathcal{M}$ 由有限生成 $A$-模 $M$ 定义，$\varphi^{*}\mathcal{M}$ 对应于 $B$-模 $B \otimes_{A} M$.$X_{\mathrm{et}}$ 上有预层 $U \mapsto \Gamma(U, \varphi^{*}\mathcal{M})$，记为 $\mathcal{M}_{\mathrm{et}}$.例如，$(\mathcal{O}_{X_{\mathrm{zar}}})_{\mathrm{et}} = \mathcal{O}_{X_{\mathrm{et}}}$.要验证 $\mathcal{M}_{\mathrm{et}}$ 是层，只需(感谢命题~\ref{prop:6.6})证明：只要 $A \to B$ 忠实平坦，序列
\[
0 \to M \to B \otimes_{A} M \rightrightarrows B \otimes_{A} B \otimes_{A} M
\]
就正合.这可以像 $M = A$ 的情形那样精确证明：再次可以假定 $A \to B$ 有截面，这允许构造收缩同伦.

\begin{example}\label{ex:6.11}
设 $X$ 是域 $k$ 上的簇.对每个态射 $\varphi \colon U \to X$，有正合序列
\[
\varphi^{*}\Omega^{1}_{X/k} \to \Omega^{1}_{U/k} \to \Omega^{1}_{U/X} \to 0
\]
(Hartshorne 1977, II.8.11).若 $\varphi$ 平展，则 $\Omega^{1}_{U/X} = 0$，从而 $\varphi^{*}\Omega^{1}_{X/k} \to \Omega^{1}_{U/k}$ 满射---由于它们是同秩的局部自由层，这蕴含\footnote{提示：利用元素在域中的，列线性无关的方阵可逆.}该映射是同构.因此，$(\Omega^{1}_{X/k})_{\mathrm{et}}$ 在 $U_{\mathrm{zar}}$ 上的限制是 $\Omega^{1}_{U/k}$.
\end{example}

\begin{exercise}\label{exc:6.12}
证明：对每个忠实平坦同态 $A \to B$ 与 $A$-模 $M$，序列
\[
0 \to M \to B \otimes_{A} M \to B \otimes_{A} B \otimes_{A} M \to B \otimes_{A} B \otimes_{A} B \otimes_{A} M \to \cdots
\]
正合.映射 $B^{\otimes i} \otimes_{A} M \to B^{\otimes (i+1)} \otimes_{A} M$ 是 $d_i \otimes 1$，其中 $d_i \colon B^{\otimes i} \to B^{\otimes (i+1)}$ 是在各个位置插入 $1$ 得到的映射的交替和.
\end{exercise}

\textbf{Spec($k$) 上的层} 设 $k$ 是域.$(\operatorname{Spec} k)_{\mathrm{et}}$ 上的 Abel 群预层 $F$ 可以看作共变函子 $\mathbf{Et}/k \to \mathbf{Ab}$(回忆 $\mathbf{Et}/k$ 是平展 $k$-代数的范畴).这样的函子是层，当且仅当对每个有限族 $(A_i)$ 的平展 $k$-代数，$F(\prod A_i) = \bigoplus F(A_i)$，并且对每个域扩张 $K/k_0$($k_0$ 在 $k$ 上有限)的有限 Galois 扩张，$F(k_0) \xrightarrow{\sim} F(K)^{\operatorname{Gal}(K/k_0)}$.

选取 $k$ 的可分闭包 $k^{\mathrm{sep}}$，设 $G = \operatorname{Gal}(k^{\mathrm{sep}}/k)$.对 $(\operatorname{Spec} k)_{\mathrm{et}}$ 上的层 $F$，定义
\[
M_F = \varinjlim F(k'),
\]
其中 $k'$ 遍历 $k^{\mathrm{sep}}$ 中在 $k$ 上有限且 Galois 的子域 $k'$.则 $M_F$ 是离散 $G$-模.

反过来，若 $M$ 是离散 $G$-模，定义
\[
F_M(A) = \operatorname{Hom}_G(F(A), M),
\]
其中 $F(A)$ 是 $G$-集合 $\operatorname{Hom}_{k\text{-alg}}(A, k^{\mathrm{sep}})$(见 \S 3).则 $F_M$ 是 $\operatorname{Spec} k$ 上的层.

函子 $F \mapsto M_F$ 与 $M \mapsto F_M$ 定义了 $(\operatorname{Spec} k)_{\mathrm{et}}$ 上层范畴与离散 $G$-模范畴之间的等价.

\subsection*{茎}

设 $X$ 是代数闭域上的簇，$F$ 是 $X_{\mathrm{et}}$ 上的预层或层.$F$ 在点 $x \in X$ 处的茎是
\[
F_{\bar{x}} = \varinjlim_{(U, u)} F(U),
\]
其中极限遍历 $x$ 的平展邻域.

设 $X$ 是概形，$F$ 是 $X_{\mathrm{et}}$ 上的预层或层.$F$ 在几何点 $\bar{x} \to X$ 处的茎是
\[
F_{\bar{x}} = \varinjlim_{(U, u)} F(U),
\]
其中极限遍历 $\bar{x}$ 的平展邻域.

\begin{example}\label{ex:6.13}
(a) $\mathcal{O}_{X_{\mathrm{et}}}$ 在 $\bar{x}$ 处的茎是 $\mathcal{O}_{X, \bar{x}}$，即 $\bar{x}$ 处的严格局部环.

(b) $F_Z$ 在 $\bar{x}$ 处的茎($Z$ 是 $X$ 上的簇或有限型概形)是 $Z(\mathcal{O}_{X, \bar{x}})$.例如，$\mu_n$,$\mathbb{G}_{a}$，$\mathbb{G}_{m}$ 与 $\mathrm{GL}_n$ 在 $\bar{x}$ 处的茎分别是 $\mu_n(\mathcal{O}_{X, \bar{x}})$,$\mathcal{O}_{X, \bar{x}}$(看作 Abel 群)，$\mathcal{O}^{\times}_{X, \bar{x}}$ 与 $\mathrm{GL}_n(\mathcal{O}_{X, \bar{x}})$.

(c) $\mathcal{M}_{\mathrm{et}}$ 在 $\bar{x}$ 处的茎($\mathcal{M}$ 是凝聚 $\mathcal{O}_X$-模)是 $M_x \otimes_{\mathcal{O}_{X, x}} \mathcal{O}_{X, \bar{x}}$，其中 $M_x$ 是 $\mathcal{M}$ 在 $x$ 处的茎(作为 Zariski 拓扑的层).

(d) 对域 $k$ 上 $\operatorname{Spec} k$ 的层 $F$，在 $\bar{x} = \operatorname{Spec} k^{\mathrm{sep}} \to \operatorname{Spec} k$ 处的茎是 $M_F$(看作 Abel 群).
\end{example}

\subsection*{摩天楼层}

一般地，层 $F$ 称为摩天楼层，如果除有限多个 $x$ 外 $F_{\bar{x}} = 0$.(回忆 $\bar{x}$ 表示 $X$ 的像为 $x \in X$ 的几何点.)我们需要一些特殊的摩天楼层.

设 $X$ 是 Hausdorff 拓扑空间，$x \in X$，$\Lambda$ 是 Abel 群.定义
\[
\Lambda_x(U) = \begin{cases}
\Lambda & \text{若 } x \in U;\\
0 & \text{其他情形}.
\end{cases}
\]
则 $\Lambda_x$ 是 $X$ 上的层.显然，$\Lambda_x$ 在 $y \neq x$ 处的茎是 0，在 $x$ 处的茎是 $\Lambda$.设 $F$ 是 $X$ 上的层.由正向极限的定义，我们看到给出同态 $F_x \overset{\mathrm{def}}{=} \varinjlim_{U} F(U) \to \Lambda$ 与给出相容的映射族 $F(U) \to \Lambda$($U$ 取遍 $x$ 的每个开邻域)是一回事，而给出这样的族与给出层映射 $F \to \Lambda_x$ 是一回事.因此
\[
\operatorname{Hom}(F, \Lambda_x) \cong \operatorname{Hom}(F_x, \Lambda).
\]

现在设 $X$ 是代数闭域上的簇，$x \in X$.对平展映射 $\varphi \colon U \to X$，定义
\[
\Lambda_x(U) = \bigoplus_{u \in \varphi^{-1}(x)} \Lambda.
\]
于是，除非 $x \in \varphi(U)$，$\Lambda_x(U) = 0$，这时它是由映到 $x$ 的 $U$ 的点指标化的 $\Lambda$ 的副本的和.同样，$\Lambda_x$ 是层，它的茎除在 $\bar{x}$ 处为零，在 $\bar{x}$ 处茎为 $\Lambda$.设 $F$ 是 $X$ 上的层，$F_{\bar{x}} \to \Lambda$ 是群同态.选取 $u \in \varphi^{-1}(x)$ 把 $(U, u)$ 实现为 $x$ 的平展邻域，从而确定映射 $F(U) \to F_{\bar{x}} \to \Lambda$.把这些映射对 $u \in \varphi^{-1}(x)$ 组合，得到同态
\[
F(U) \to \Lambda_x(U) \overset{\mathrm{def}}{=} \bigoplus_{u} \Lambda.
\]
它们与限制映射相容，从而定义同态 $F \to \Lambda_x$.这样，我们再次得到同构
\[
\operatorname{Hom}(F, \Lambda_x) \cong \operatorname{Hom}(F_{\bar{x}}, \Lambda).
\]

设 $X$ 是概形，$i \colon \bar{x} \to X$ 是 $X$ 的几何点.对任何平展映射 $\varphi \colon U \to X$，定义
\[
\Lambda_{\bar{x}}(U) = \bigoplus_{\operatorname{Hom}_X(\bar{x}, U)} \Lambda.
\]
同样，这给出 $X_{\mathrm{et}}$ 上的层，并有自然同构 $\operatorname{Hom}(F, \Lambda_{\bar{x}}) \to \operatorname{Hom}(F_{\bar{x}}, \Lambda)$.然而，若 $x \overset{\mathrm{def}}{=} i(\bar{x})$ 不闭，则当 $\bar{y}$ 的像不等于 $x$ 时，$(\Lambda_{\bar{x}})_{\bar{y}} = 0$ 未必成立.因此，除非 $x$ 闭，$\Lambda_{\bar{x}}$ 未必是摩天楼层.

\begin{example}\label{ex:6.14}
设 $X = \operatorname{Spec} k$，几何点 $\bar{x} \to X$ 对应于 $k$ 到其可分闭包 $k^{\mathrm{sep}}$ 的包含.则 $\Lambda_{\bar{x}}$ 是 $X_{\mathrm{et}}$ 上对应于如下离散 $G$-模的层：该模由连续映射 $G \to \Lambda$ 组成(诱导 $G$-模---见 Serre, \emph{Cohomologie Galoisienne}, 2.5).
\end{example}

\subsection*{局部常值层}

设 $X$ 是拓扑空间.$X$ 上的层 $F$ 称为局部常值的，如果对 $X$ 的某个开覆盖中的所有 $U$，$F|_{U}$ 都是常值的.

当 $X$ 连通，局部道路连通且局部单连通时，局部常值层由 $X$ 的基本群分类.

固定点 $x \in X$，设 $F$ 是 $X$ 上的局部常值层.设 $\gamma \colon [0, 1] \to X$ 是满足 $\gamma(0) = x = \gamma(1)$ 的连续映射，即 $\gamma$ 是 $X$ 中基于 $x$ 的闭路.$F$ 在 $[0, 1]$ 上的原像 $\gamma^{*}(F)$ 是常值的.选取 $\gamma^{*}(F)$ 与常值层 $\Lambda$ 的同构，确定了对每个 $0 \leq a \leq 1$ 的同构 $F_{\gamma(a)} \to \Lambda$.由于 $\gamma(0) = \gamma(1) = x$，得到两个同构 $F_x \to \Lambda$，它们相差 $F_x$ 的一个自同构 $\alpha(\gamma)$.映射 $\gamma \mapsto \alpha(\gamma)$ 定义同态 $\pi_1(X, x) \to \operatorname{Aut}(F_x)$.下面是众所周知的：

\begin{proposition}\label{prop:6.15}
映射 $F \mapsto F_x$ 定义了 $X$ 上集合(分别地，Abel 群)的局部常值层范畴与 $\pi_1(X, x)$ 集合(分别地，模)范畴之间的等价.
\end{proposition}

设 $X$ 是代数簇(或概形).$X_{\mathrm{et}}$ 上的层 $F$ 是局部常值的，如果对 $X$ 的某个平展覆盖中的所有 $U \to X$，$F|_{U}$ 都是常值的.

\begin{proposition}\label{prop:6.16}
假定 $X$ 连通，设 $\bar{x}$ 是 $X$ 的几何点.映射 $F \mapsto F_{\bar{x}}$ 定义了 $X$ 上带有限茎的集合(分别地，Abel 群)的局部常值层范畴与有限 $\pi_1(X, \bar{x})$ 集合(分别地，模)范畴之间的等价.
\end{proposition}

\begin{proof}
事实上，可以证明：若 $Z \to X$ 是有限平展映射，则层 $F_Z$ 局部常值且茎有限，而每个这样的层对某个 $Z$ 都是这种形式的；更精确地说，$Z \mapsto F_Z$ 定义了从 $\mathbf{FEt}/X$ 到 $X_{\mathrm{et}}$ 上带有限茎的局部常值层范畴的等价.现在可以应用定理~\ref{thm:3.1} 证明命题.

我概述上述论断的一个证明.考虑满平展态射 $Z \to X$.如上面 6.7 所述，存在满有限平展映射 $Z' \to Z$，使得复合 $Z' \to X$ 是 Galois 覆盖.可以证明 $Z \times_{X} Z'$ 是 $Z'$ 的副本的不交并.层 $F_Z$ 限制为 $Z'$ 上的层 $F_{Z \times_{X} Z'}$，它是常值的.这表明 $F_Z$ 局部常值，并且事实上在有限平展覆盖 $Z' \to Z$ 上变成常值的.

反过来，设 $F$ 是带有限纤维的局部常值层.按假设，存在平展覆盖 $(U_i \to X)_{i \in I}$，使得 $F|_{U_i}$ 常值，从而具有形式 $F_{Z_i}$，其中 $Z_i$ 在 $U_i$ 上有限且平展(事实上，$Z_i$ 是 $U_i$ 的副本的不交并).同构 $(F|_{U_i})|_{U_i \times_{X} U_j} \xrightarrow{\sim} (F|_{U_j})|_{U_i \times_{X} U_j}$ 诱导同构 $Z_i \times_{U_i} (U_i \times_{X} U_j) \to Z_j \times_{U_j} (U_i \times_{X} U_j)$，而下降理论(比较 EC I 2.23)表明由诸 $Z_i$ 与这些同构组成的系统来自某个 $Z \to X$.
\end{proof}
